\documentclass[12pt,a4paper]{article}
\usepackage{fontspec}
\newfontfamily\cjkfont{Droid Sans Fallback}
\newcommand{\zh}[1]{{\cjkfont #1}}
\usepackage[spanish]{babel}
\usepackage{geometry}
\geometry{margin=1in}
\usepackage{hyperref}
\usepackage{setspace}
\onehalfspacing
\usepackage{parskip}

\begin{document}

% --- Cover page ---
\begin{titlepage}
\centering
\vspace*{3cm}

{\LARGE\bfseries Problemas y políticas públicas en China\par}

\vspace{2cm}

{\large Hesus García Cobos\par}

\vspace{1cm}

{\large Profesora: Socorro Guevara Salazar\par}

\vspace{1cm}

{\large 20 de marzo de 2026\par}

\vfill
\end{titlepage}

% --- Body ---

\section*{Dos problemas relevantes en China en el área de educación e inteligencia artificial}

\subsection*{1. La brecha digital entre ciudades y zonas rurales}

China tiene más de 200 millones de estudiantes en educación básica. El problema es que las condiciones para enseñar IA no son las mismas en Shanghái que en una escuela rural de Guizhou o Gansu. La propia guía del Ministerio de Educación de 2024 reconoce esta brecha y pide ``desarrollo equilibrado entre zonas urbanas y rurales'' y ``mayor apoyo para escuelas en áreas rurales y remotas'' (\zh{教育部}, 2024). En la práctica, muchas escuelas rurales no tienen laboratorios, conectividad estable ni maestros capacitados en IA. El gobierno designó 184 escuelas piloto como bases de educación en IA, pero la mayoría están en ciudades grandes. Hay un programa de ``asistencia pareada'' (\zh{结对帮扶}) entre escuelas urbanas y rurales, pero la distancia entre las condiciones de unas y otras sigue siendo enorme.

\subsection*{2. Falta de maestros capacitados para enseñar IA}

El segundo problema es quién enseña. Desde septiembre de 2025, todas las escuelas primarias y secundarias deben ofrecer mínimo 8 horas anuales de IA. Pero no hay suficientes maestros preparados para cubrir esa demanda en un país de ese tamaño. La guía de 2024 propone reclutar profesionales de universidades y empresas de tecnología como instructores de medio tiempo (\zh{教育部}, 2024). Que un gobierno tenga que recurrir a profesionales externos para cubrir sus aulas indica que el sistema de formación docente no alcanza. Y un profesional que sabe de IA no necesariamente sabe enseñar, que es otro problema aparte.

\section*{Resumen de dos políticas públicas chinas relevantes}

\subsection*{1. Plan de desarrollo de inteligencia artificial de nueva generación (2017)}

El Consejo de Estado publicó este plan en julio de 2017 con la meta de convertir a China en líder mundial de IA para 2030. Establece tres etapas: alcanzar el nivel de los países avanzados para 2020, lograr avances importantes para 2025, y liderar a nivel mundial para 2030. En educación, el plan pide construir un sistema educativo en IA que incluya cursos desde primaria y formación de talento en todos los niveles. Es el documento que le dio a todo el sistema educativo chino la instrucción de moverse hacia la inteligencia artificial (\zh{国务院}, 2017).

\subsection*{2. Guía para la educación en IA en primaria y secundaria (2024)}

En diciembre de 2024, el Ministerio de Educación emitió esta guía que organiza la enseñanza de IA por niveles: percepción y experiencia en los primeros grados de primaria, comprensión y aplicación en grados superiores de primaria y secundaria, y creación de proyectos en preparatoria. Establece el uso de la plataforma nacional de educación inteligente, la creación de bases de educación en IA por lotes, y el intercambio de recursos entre escuelas urbanas y rurales. Es la política que concretó las metas del plan de 2017 en un currículo real (\zh{教育部}, 2024).

% --- References ---
\newpage
\begin{center}
{\large\bfseries Referencias}
\end{center}
\vspace{0.5cm}

\noindent \zh{国务院} [Consejo de Estado]. (2017). \zh{国务院关于印发新一代人工智能发展规划的通知} [Plan de desarrollo de inteligencia artificial de nueva generación]. \url{https://www.gov.cn/xinwen/2017-07/20/content_5212064.htm}

\vspace{0.4cm}
\noindent \zh{教育部} [Ministerio de Educación]. (2024). \zh{教育部办公厅关于加强中小学人工智能教育的通知} [Guía para la educación en IA en primaria y secundaria]. \url{https://www.moe.gov.cn/srcsite/A06/s3321/202412/t20241202_1165500.html}

\end{document}
